\section{Два общих препятствия для семейства $H(3,3n+1,3n+2)$}

Работаем над $\mathbb C$; результаты тем самым применимы над любым
подполем $\mathbb C$. Используется нормировка
$C_j(t)=2T_j(t/2)$ и обозначение $U_j=U_j(x/2)$ для многочленов
Чебышёва второго рода; $U_{-1}=0$, $U_0=1$.
Для $n\geq1$ положим
\[
 \gamma_n(t)=(C_3(t),C_{3n+1}(t),C_{3n+2}(t)).
\]
Три индекса попарно взаимно просты. Полученные ниже препятствия
не доказывают невыпрямляемость этих вложений.

\subsection{Запрет координат, линейных по одной поперечной переменной}

\paragraph{Теорема.}
Для любого $n\geq1$ не существует $H\in\mathbb C[x,y,z]$, одновременно
удовлетворяющего условиям
\[
 H\circ\gamma_n=t,\qquad
 (x,H,W)\in\operatorname{Aut}_{\mathbb C}\mathbb A^3
 \quad\text{для некоторого }W,
\]
и имеющего хотя бы одну из форм
\[
 H=Q(x,y)+M(x)z,\qquad H=Q(x,z)+M(x)y,
\]
где $Q,M$ --- произвольные многочлены без ограничений степени.
Заключение сохраняется при замене $t$ в первом условии на $at+b$,
$a\ne0$, посредством аффинной нормировки $H$.

\paragraph{Доказательство.}
В кольце $\mathbb C[x,t]/(t^3-3t-x)$ положим $s=t^2-2$.
Это свободный модуль с базисом $1,t,s$. Положим
$u=U_n$ и $v=U_{n-1}$. Рекурсия Чебышёва даёт
\[
 y=ut-vs,\qquad z=-vt+us,\qquad
 u^2+v^2-xuv=1.
\]
Поэтому
\[
 y^2=2+(-2uv+xv^2)t+(u^2-v^2)s,
\]
\[
 z^2=2+(-2uv+xu^2)t-(u^2-v^2)s.
\]
Используем также тождества
\[
 u^2-v^2=U_{2n},
\]
\[
 u^3-3uv^2+xv^3=U_{3n},\qquad
 v^3-3vu^2+xu^3=U_{3n+1}.
\]
Все эти тождества следуют либо из рекурсии, либо из
$U_j(\cos\theta)=\sin((j+1)\theta)/\sin\theta$.

Рассмотрим первую форму. Моническое соотношение
\[
 y^3-3y-C_{3n+1}(x)=0
\]
позволяет заменить $Q(x,y)$ при ограничении на кривую остатком
$A(x)y^2+B(x)y+C(x)$. Сравнение коэффициентов $t,s$ и исключение
$B$ даёт
\[
 \boxed{U_{3n}A+U_{2n}M=U_{n-1}.} \tag{I}
\]
Если $M=0$, это невозможно по степеням. Пусть $M\ne0$ и
$M(\alpha)=0$. Редукция автоморфизма $(x,H,W)$ при $x=\alpha$
показывает, что $Q(\alpha,y)$ --- координата $\mathbb C[y,z]$,
зависящая только от $y$. Следовательно, она аффинна по $y$.
Деление на монический многочлен коммутирует со специализацией,
поэтому $A(\alpha)=0$. Из (I) следует $U_{n-1}(\alpha)=0$.
Итак, все корни $M$ лежат среди вещественных корней $U_{n-1}$.

Два наибольших корня $U_{3n}$ равны
\[
 x_j=2\cos\frac{j\pi}{3n+1},\qquad j=1,2.
\]
При $n\geq2$ оба лежат строго правее наибольшего корня
$2\cos(\pi/n)$ многочлена $U_{n-1}$.
При $n=1$ многочлен $U_{n-1}=1$ не имеет корней, и $M$ постоянен.
Для $\theta_j=j\pi/(3n+1)$ имеем
\[
 \left.\frac{U_{n-1}}{U_{2n}}\right|_{x=x_j}
 =\frac{\sin(n\theta_j)}{\sin((2n+1)\theta_j)}
 =(-1)^{j+1}.
\]
Числитель и знаменатель ненулевые. Следовательно,
\[
 M(x_1)=1,\qquad M(x_2)=-1.
\]
Но $M=c\prod_\alpha(x-\alpha)^{e_\alpha}$, все $\alpha$ вещественны
и меньше $x_2$. Поэтому $M(x_1)/M(x_2)$ --- положительное
вещественное число, даже если $c$ комплексно. Противоречие.

Для второй формы редуцируем $Q(x,z)$ по
$z^3-3z-C_{3n+2}(x)$ до $Az^2+Bz+C$. Тот же расчёт даёт
\[
 \boxed{U_{3n+1}A+U_{2n}M=U_n.} \tag{II}
\]
Все корни $M$ должны быть корнями $U_n$. Два наибольших корня
\[
 x_j=2\cos\frac{j\pi}{3n+2},\qquad j=1,2,
\]
многочлена $U_{3n+1}$ лежат правее всех корней $U_n$, поскольку
$2/(3n+2)<1/(n+1)$ для $n\geq1$. На этих корнях
\[
 \left.\frac{U_n}{U_{2n}}\right|_{x=x_j}
 =\frac{\sin((n+1)\theta_j)}{\sin((2n+1)\theta_j)}
 =(-1)^{j+1},\qquad \theta_j=\frac{j\pi}{3n+2}.
\]
Опять получаем $M(x_1)=1$, $M(x_2)=-1$ на луче справа от всех
корней $M$, что невозможно.\hfill$\square$

\paragraph{Почему это не противоречит известным выпрямлениям.}
Уже $n=1$ даёт $H(3,4,5)$, для которого в репозитории имеется
проверяемый выпрямитель над $\mathbb Q(\sqrt{-3})$.
Его параметрическая компонента нелинейно зависит от обеих
поперечных переменных; она не принадлежит ни одной запрещённой форме.
Теорема показывает необходимость выхода из двух простых классов,
а не невозможность такого выхода.

\paragraph{Следствие об абсолютных координатах.}
В формулировке теоремы можно заменить требование
$(x,H,W)\in\operatorname{Aut}\mathbb A^3$ более слабым требованием:
$H$ является какой-либо координатой $\mathbb C[x,y,z]$,
без условия сохранения $x$.

Действительно, $M=0$ уже исключается тождеством (I) или (II).
При $M\ne0$ выберем $\alpha$ с $M(\alpha)\ne0$.
Тогда $H(\alpha,y,z)$ --- координата $\mathbb C[y,z]$,
поскольку линейна по одной переменной с ненулевым постоянным коэффициентом.
Абсолютная координатность $H$ даёт
$\mathbb C[x,y,z]/(H)\simeq\mathbb C^{[2]}$.
По Proposition~3.1 статьи Edo--van den Essen--Maubach,
\emph{A note on $k[z]$-automorphisms in two variables},
\url{https://arxiv.org/abs/0809.0767},
после переименования фиксированной переменной $z$ в $x$ эти два
условия влекут, что $H$ является $\mathbb C[x]$-координатой.
Это противоречит доказанной теореме.\hfill$\square$

Основное относительное препятствие доказано выше элементарно.
Это усиление использует указанную классическую теорему;
оно не объявляется следствием одной проверки якобиана.
Таким образом, у любого выпрямляющего автоморфизма модели $\gamma_n$
компонента, восстанавливающая параметр, должна выходить из обеих
указанных форм даже тогда, когда другие компоненты меняют $x$.

\subsection{Запрет естественного переноса через $C_n(x)=0$}

\paragraph{Теорема о центре переноса.}
Для $n\geq1$ рассмотрим
\[
 \beta_n(x,y,z)=(x,C_n(x)z-y^2+2,y).
\]
Тогда
\[
 \beta_n\circ\gamma_n=(C_3,C_2,C_{3n+1}).
\]
Не существует многочлена $H(x,Y,Z)$, который одновременно
удовлетворял бы
\[
 H(C_3(t),C_2(t),C_{3n+1}(t))=t
\]
и имел аффинный центр
$H(\alpha,2-y^2,y)$ по $y$ на каком-либо корне
$\alpha$ многочлена $C_n(x)$.
Поэтому этот перенос не проходит критерий относительной координаты
ни при каком выборе параметрического $H$ и ни при какой его степени.

\paragraph{Доказательство.}
Тождество
\[
 C_{3n}C_{3n+2}-C_{3n+1}^2+2=C_2
\]
следует из $C_aC_b=C_{a+b}+C_{|a-b|}$ и
$C_{3n}=C_n(C_3)$.
Пусть $C_n(\alpha)=0$. В кольце
$\mathbb C[t]/(t^3-3t-\alpha)$ положим
$L=C_{3n+1}(t)$. Тогда $C_2(t)=2-L^2$.
Предположив $H(\alpha,2-y^2,y)=\lambda y+\mu$, получим
\[
 t=\lambda L+\mu
 =\lambda\left(U_n(\alpha/2)t-
 U_{n-1}(\alpha/2)(t^2-2)\right)+\mu.
\]
Но $U_{n-1}(\alpha/2)\ne0$. Это следует из тождества
\[
 C_n(x)^2+(4-x^2)U_{n-1}(x/2)^2=4.
\]
Сравнение коэффициентов в свободном базисе $1,t,t^2$ даёт
сначала $\lambda=0$, затем противоречие коэффициенту $t$.
Аргумент не требует различности корней $t^3-3t-\alpha$.
\hfill$\square$

При нечётном $n$ целевая тройка вообще не является замкнутым
вложением: $3n+1$ чётно, и $t=\sqrt3$ и $t=-\sqrt3$ имеют
одинаковые три координаты. При чётном $n$ целевая тройка
$(3,2,3n+1)$ является выпрямляемым гармоническим вложением,
но сформулированное препятствие центру сохраняется.

Последнее выпрямление можно предъявить непосредственно.
Если $n=2m$, то $C_{6m+1}(t)=tP_m(t^2-2)$ для некоторого
$P_m\in\mathbb Q[Y]$, причём
$P_m(1)=(-1)^m=:\varepsilon$.
Положим
\[
 A_m(Y)=\frac{1-\varepsilon P_m(Y)}{Y-1}\in\mathbb Q[Y].
\]
Тогда
\[
 H_0(x,Y,Z)=xA_m(Y)+\varepsilon Z,
 \qquad H_0(C_3,C_2,C_{6m+1})=t.
\]
Поскольку $\varepsilon=\pm1$, отображение $(x,H_0,Y)$
есть полиномиальный автоморфизм. Следовательно, центр не удаётся
исправить даже несмотря на наличие явной относительной координаты
на целевом вложении.

\subsection{Статус и воспроизводимость}
Оба доказательства относятся к бесконечному семейству и не опираются
на ограниченный перебор. Скрипт \texttt{verify\_uniform\_obstructions.py}
проверяет 160 точных полиномиальных тождеств для $1\leq n\leq20$.
Конечный прогон служит проверкой реализации, а не доказательством
утверждений для всех $n$.
Мотивы, $K$-теория и гомологическая алгебра не используются.
Литературная новизна этих следствий не утверждается.
